\section{Case Study: Citation Verification}
\label{sec:casestudy}

The six families of Section~\ref{sec:setup} are controlled by construction. To
check that the framework says something useful about a task people actually
deploy language models on, we apply it to reference generation---a setting
where fabricated output is a known and consequential failure mode, and where an
external registry supplies a near-perfect verifier for the price of an HTTP
request.

The task ``produce a real published reference on topic $X$'' maps onto the
framework exactly. A generator proposes $k$ candidate references; a selector
picks one; ground truth is whether the reference exists. We ran this with
Qwen2.5-1.5B over \nCitTopics{} research topics spanning the sciences, drawing
$k=8$ candidate references per topic, and resolved every candidate against the
OpenAlex catalogue by fuzzy title match (threshold $0.88$), which supplies both
ground truth and one of the two selectors under comparison.

\paragraph{Resolution failures are not fabrications.}
A registry lookup can fail because the paper does not exist or because the
service is temporarily unavailable, and conflating the two would convert a
network condition into a measured hallucination. We retry with exponential
backoff and record any candidate we still cannot resolve as \emph{unresolved}
rather than absent; \nCitUnresolved{} candidates were excluded on this basis,
and a topic contributes to the analysis only if all eight of its candidates
resolved.

\subsection{Results}

The generator hallucinates heavily: \citHallucRate{} of proposed references
could not be found, so a single candidate is real with probability
$\citPone{}$. Coverage is much higher---at least one of eight candidates
resolves for $\citCoverage{}$ of topics---so the headroom is large and the
entire question is which selector can convert it.

\begin{description}[leftmargin=0pt,itemsep=1pt,topsep=2pt,font=\normalfont\itshape]
\item[The model as its own critic.] Asked whether each reference is real, the
  model answers ``real'' for \citSaysReal{} of candidates, including most of the
  fabricated ones. Its pairwise gap on this task is $\gap = \citG{}$, and its
  realised selection efficiency is $\eff = \citEtaLLM{}$, at a cost of
  \citLLMTokens{} tokens.
\item[The external registry.] The same selection made by an OpenAlex lookup
  attains $\eff = \citEtaAPI{}$ at zero model tokens, raising accuracy from
  $\citPone{}$ to $\citAccAPI{}$.
\end{description}

The latent-score model of Proposition~\ref{prop:latent}, given only the measured
$\gap$, predicts $\hat{\eff} = \citPredEta{}$ for the critic---derived without
running the selection at all.

\subsection{What the case study shows}

Two agents that look alike in an architecture diagram---one checks citations,
one polishes prose---are not alike at all under this framework. The first has an
external oracle available for the price of an HTTP request and converts nearly
all of its coverage headroom. The second has no oracle, and a model asked to
supply one performs near chance. The generalisable lesson is that the question
worth asking about a proposed agent is not how capable it is but whether
anything cheap can tell when it is wrong. Where an external check exists, use
it: in our measurements it dominated the model-critic alternative on both
accuracy and cost, which is not a trade-off but a strict improvement.
